%----------HEADING-----------------
\begin{tabular*}{\textwidth}{l@{\extracolsep{\fill}}r}
  \textbf{\href{https://www.linkedin.com/in/yufeng-xia-12648a25b}{\Large Yufeng (Alex) Xia}} & Email: \href{mailto:yufeng.xia@ed.ac.uk}{yufeng.xia@ed.ac.uk} \\
  \href{https://www.linkedin.com/in/yufeng-xia-12648a25b}{linkedin.com/in/yufeng-xia-12648a25b} & Mobile: \href{tel:+447918155093}{(+44) 7918155093} \\
  Edinburgh, UK & \href{https://xyf2002.github.io/}{xyf2002.github.io} \\
\end{tabular*}

\vspace{-2pt}

%-----------PROFILE-----------------
\section{Profile}
  \begin{itemize}[leftmargin=*]
    \item\small{MSc by Research student at the Institute for Computing Systems Architecture (ICSA), University of Edinburgh, supervised by Prof.\ Mahesh K.\ Marina. My interest lies at the intersection of \textbf{distributed systems, machine learning systems, and performance engineering}. My work spans \textbf{accelerator sharing and scheduling}, where stateful ML workloads run on contended, reclaimable GPU capacity alongside a latency-critical tenant, as well as \textbf{large-scale emulation and profiling}, the infrastructure used to measure how such systems behave at cluster scale.
    \par\vspace{3pt}
    I work on the layer that decides which workload gets which accelerator, and on what happens when that decision changes: preemption tolerance, state preservation, and graded degradation under resource pressure. These systems have led to publications at top-tier venues including ACM MobiCom, and what interests me most now is applying the same scheduling and state-management questions to large-scale AI inference infrastructure. \textbf{Available full time from September 2026.}}
  \end{itemize}
